\section{Serving Economics}\label{sec:econ}

Claim~A is an economic claim, and it cannot rest on architecture alone: a hybrid that decodes faster but costs more per task has not solved the deployment problem.
% 中文：主张 A 是经济主张，不能只靠架构论证：解码更快但每任务更贵的混合模型并未解决部署问题。
This section fixes the accounting before any number is measured, so that the measurement cannot be chosen after the fact.
% 中文：本节先固定记账口径，再谈测量，避免「事后选口径」。

\paragraph{Accounting.}
For one benchmark task with $R$ \texttt{run\_code} rounds, mean context $L$ and $G_{\mathrm{out}}$ generated tokens, the local cost is
% 中文：记账。对一个含 R 轮 \texttt{run\_code}、平均上下文 L、生成 $G_{out}$ 个 token 的基准任务，本地成本为
\begin{equation}
C_{\text{task}} \;=\; \underbrace{\frac{P_{\text{pre}}\,L\,R}{S_{\text{pre}}}}_{\text{prefill GPU-s}}
\;+\; \underbrace{\frac{G_{\mathrm{out}}}{S_{\text{dec}}}}_{\text{decode GPU-s}}
\;+\; \underbrace{\frac{\mu\,M_{\max}\,T_{\text{wall}}}{M_{\text{gpu}}}}_{\text{memory occupancy}},
\label{eq:cost}
\end{equation}
where $S_{\text{pre}},S_{\text{dec}}$ are measured throughputs, $M_{\max}$ the peak session state of Eq.~\ref{eq:mem}, and $\mu$ the amortised accelerator price; the API comparison uses the provider's per-token price at the same token counts.
% 中文：其中 $S_{pre},S_{dec}$ 为实测吞吐，$M_{max}$ 为式 \ref{eq:mem} 的峰值会话状态，$\mu$ 为摊销的加速器价格；API 对照用同 token 量下的供应商按 token 计价。
Memory occupancy is kept as a separate term precisely because it is the term a hybrid improves and a throughput-only comparison hides.
% 中文：把内存占用单列，正是因为它才是混合模型改善的那一项，而只比吞吐会把它藏起来。

\begin{table}[t]
\centering
\caption{Serving economics per task (T15). Local rows use Eq.~\ref{eq:cost}; the API row is priced at the provider's rate.}
\label{tab:econ}
\scriptsize
\setlength{\tabcolsep}{2.8pt}
\begin{tabular}{@{}lccccc@{}}
\toprule
\textbf{Configuration} & \textbf{Prefill} & \textbf{Decode} & \textbf{Mem} & \textbf{GPU-s} & \textbf{Rel.\ cost} \\
 & (tok/s) & (tok/s) & (GB) & /task & \\
\midrule
\llm-2.2B (hybrid, local) & \tbd & \tbd & \tbd & \tbd & $1.0\times$ \\
dense-2.5B (local) & \tbd & \tbd & \tbd & \tbd & \tbd \\
API (ZhuLong, V4-Pro) & --- & --- & --- & \tbd & \tbd \\
\bottomrule
\end{tabular}
\tnote{Throughputs must be measured on the production stack; the Megatron-Core direct-drive numbers used during pre-training (prefill $23{,}692$ vs.\ $21{,}436$ tok/s; decode $22.77$ vs.\ $2.21$ tok/s) are \emph{not} used here because the dense decode figure is CPU-launch-bound and would overstate the gap. \rv{Reverse contingency: if memory occupancy is negligible at the measured session lengths, Claim~A must be stated on prefill/decode cost alone, and the architecture argument is reported as unrealised at benchmark scale.}}
\end{table}

\begin{figure}[t]
\centering
\phfig{3.0cm}{\textbf{Figure placeholder: cost--accuracy plane}\\[2pt]\llm vs.\ six ZhuLong backbones\\x: GPU-s per task (log); y: pooled Pass@1;\\dashed: $53.2\%$ floor of Claim~A}
\caption{Cost--accuracy plane on which Claim~A is decided. ZhuLong backbone points are placed from Table~\ref{tab:backbone} cost columns where available; the \llm point is the measurement of Table~\ref{tab:econ}.}
\label{fig:pareto}
\end{figure}

Two caveats are part of the claim, not footnotes.
% 中文：两点保留意见属于主张本身，而非脚注。
First, an API deployment carries no memory term but carries the provider's margin, so the comparison must price both at equal token counts rather than at list price.
% 中文：其一，API 部署没有内存项，却有供应商利润，因此比较须在相同 token 量下计价，而非按标价。
Second, the local model's advantage is bounded by utilisation: a single session on one GPU looks expensive until the same weights serve a team.
% 中文：其二，本地模型优势受利用率约束：单会话独占一张卡显得昂贵，直到同一份权重服务整个团队。
We therefore report both $N{=}1$ and an $N{=}8$ multiplexed case, and mark the crossover as measured or unmeasured rather than asserted.
% 中文：因此我们同时报告 N=1 与 N=8 复用两种情形，并把交叉点标注为「已测/未测」，而非断言。
